\chapter{Resultados Fundamentais da Teoria dos Números}

\section{Algoritmo de Euclides}

O Algoritmo de Euclides é um método eficiente para calcular o maior divisor
comum ($\gcd(a, b)$) de dois inteiros $a$ e $b$, não ambos nulos. Baseado na
propriedade $\gcd(a, b) = \gcd(b, a \mod b)$, usa divisões sucessivas até
alcançar um resto zero. Sua relevância é fundamental em teoria dos números,
simplificação de frações, e como base para algoritmos avançados, como o
Euclides Estendido e aplicações em criptografia.

A Proposição da Propriedade Fundamental do Algoritmo de Euclides estabelece que
o maior divisor comum de dois inteiros $a$ e $b$ é igual ao maior divisor
comum de $b$ e o resto da divisão de $a$ por $b$. Essa propriedade é a
pedra angular do Algoritmo de Euclides, permitindo calcular o $\gcd(a, b)$
por divisões sucessivas. Sua relevância é central em teoria dos números,
sustentando aplicações como simplificação de frações, criptografia e a
Identidade de Bézout.

\begin{proposition}[Propriedade Fundamental do Algoritmo de Euclides]
Para quaisquer inteiros $a$ e $b$, com $b \neq 0$, temos:
\[
\gcd(a, b) = \gcd(b, a \mod b).
\]
\end{proposition}

\begin{proof}
Seja $a \mod b = r$, onde $a = q b + r$ e $0 \leq r < |b|$, com $q = \lfloor a / b \rfloor$. Precisamos mostrar que $\gcd(a, b) = \gcd(b, r)$.

\begin{enumerate}[noitemsep]
    \item Primeiro vamos mostrar que se $d$ divide $a$ e $b$, então $d$ divide $b$ e $r$.
          Suponha $d \mid a$ e $d \mid b$, ou seja, $a = k d$ e $b = m d$ para inteiros $k, m$.
          Como $r = a - q b$, temos: $r = k d - q \cdot m d = (k - q m) d$.
          Logo, $d \mid r$, pois $k - q m$ é inteiro.
          Assim, $d$ é divisor comum de $b$ e $r$.

    \item Agora, mostraremos que se $d$ divide $b$ e $r$, então $d$ divide $a$ e $b$.
          Suponha $d \mid b$ e $d \mid r$, ou seja, $b = m d$ e $r = n d$ para inteiros $m, n$.
          Como $a = q b + r$, temos: $a = q \cdot m d + n d = (q m + n) d$.
          Logo, $d \mid a$, pois $q m + n$ é inteiro.
          Assim, $d$ é divisor comum de $a$ e $b$.
    \item Por fim,
          os divisores comuns de $a$ e $b$ são os mesmos que os de $b$ e $r$.
          Desta forma, $\gcd(a, b) = \gcd(b, r) = \gcd(b, a \mod b)$.
\end{enumerate}

Caso especial: Se $b = 0$, $\gcd(a, 0) = |a|$, e a propriedade não se
  aplica diretamente, pois $a \mod 0$ não é definido. Nesse caso, o Algoritmo
  de Euclides termina com $\gcd(a, 0) = |a|$.
\end{proof}



\begin{algorithm}
\caption{Algoritmo de Euclides}
\begin{algorithmic}[1]
\Require Inteiros $a$, $b$, não ambos nulos
\Ensure $\gcd(a, b)$
\State $a \gets |a|$, $b \gets |b|$ \Comment Trata valores negativos
\While{$b \neq 0$}
    \State $r \gets a \mod b$
    \State $a \gets b$
    \State $b \gets r$
\EndWhile
\Return $a$
\end{algorithmic}
\end{algorithm}




\section{Algoritmo de Euclides Estendido}

O Algoritmo de Euclides Estendido além de calcular o maior divisor comum
($\gcd(a, b)$) de dois inteiros $a$ e $b$,
também encontra os coeficientes inteiros $m$ e $n$ tais que $am + bn = \gcd(a, b)$.
Enquanto o algoritmo padrão foca apenas no $\gcd$, o estendido rastreia
combinações lineares dos números iniciais. Sua relevância é ampla: é essencial
na resolução de equações diofantinas, na criptografia (como no RSA para
inversos modulares) e na simplificação de frações, sendo uma ferramenta chave
em teoria dos números.

O Algoritmo de Euclides Estendido mantém a mesma sequência de divisões do
Algoritmo de Euclides, mas adiciona:
\begin{itemize}[noitemsep]
    \item Rastreamento de coeficientes: Inicializa $x[0] = 1$, $y[0] = 0$,
      $x[1] = 0$, $y[1] = 1$ e atualiza $x[n+1] = x[n-1] - q \cdot x[n]$,
      $y[n+1] = y[n-1] - q \cdot y[n]$ a cada iteração, onde $q$ é o quociente
      da divisão.
    \item Saída adicional: Além do $\gcd(a, b)$, retorna $m = x[n]$ e $n =
      y[n]$, que satisfazem $am + bn = \gcd(a, b)$.
\end{itemize}
Veja sua formulação no \Cref{alg:euclides}.

\begin{algorithm}
\caption{Algoritmo de Euclides Estendido}\label{alg:euclides}
\begin{algorithmic}[1]
\Require Inteiros $a$, $b$, não ambos nulos
\Ensure $\gcd(a, b)$ e inteiros $m$, $n$ tais que $am + bn = \gcd(a, b)$
\State $a \gets |a|$, $b \gets |b|$ \Comment Trata valores negativos
\State $x[0] \gets 1$, $y[0] \gets 0$
\State $x[1] \gets 0$, $y[1] \gets 1$
\State $n \gets 1$
\While{$b \neq 0$}
    \State $q \gets \lfloor a / b \rfloor$, $r \gets a - q \cdot b$
    \State $a \gets b$, $b \gets r$
    \State $x[n+1] \gets x[n-1] - q \cdot x[n]$
    \State $y[n+1] \gets y[n-1] - q \cdot y[n]$
    \State $n \gets n + 1$
\EndWhile
\Return $a$, $x[n]$, $y[n]$ \Comment $\gcd(a, b)$, $m$, $n$
\end{algorithmic}
\end{algorithm}





\section{Identidade de Bézout}

A Identidade de Bézout afirma que, para quaisquer inteiros $a$ e $b$, não
ambos nulos, existem inteiros $m$ e $n$ tais que $am + bn = \gcd(a, b)$,
onde $\gcd(a, b)$ é o maior divisor comum. Comprovada pelo Algoritmo de
Euclides Estendido, ela expressa o $\gcd$ como uma combinação linear de $a$
e $b$. Sua relevância é crucial em teoria dos números, permitindo resolver
equações diofantinas lineares, determinar inversos modulares em criptografia e
estudar propriedades de divisibilidade, com aplicações em matemática pura e
computacional.


\begin{theorem}[Identidade de Bézout]\label{thm:bezout}

  Sejam $a$ e $b$ inteiros não ambos zero. Então, o maior divisor comum de $a$
  e $b$, denotado $\text{gcd}(a, b) = d$, pode ser expresso como uma combinação
  linear de $a$ e $b$, ou seja, existem inteiros $x$ e $y$ tais que:

  \begin{equation}
  ax + by = d
  \end{equation}

  Além disso, $d$ é o menor inteiro positivo que pode ser escrito como uma
  combinação linear de $a$ e $b$, e qualquer inteiro que seja uma
  combinação linear de $a$ e $b$ é um múltiplo de $d$.
\end{theorem}

\begin{proof}
  Seja $d = \text{gcd}(a, b)$. Vamos provar que existem inteiros $x$ e $y$ tais
  que $ax + by = d$, e que $d$ é o menor inteiro positivo expresso como
  combinação linear de $a$ e $b$.

  Primeiro, mostramos que $d$ pode ser escrito como uma combinação linear de
  $a$ e $b$. Considere o conjunto $S = \{ ax + by \mid x, y \in \mathbb{Z}, ax + by > 0 \}$,
  que contém todos os inteiros positivos que podem ser escritos
  como combinações lineares de $a$ e $b$. Como $a$ e $b$ não são ambos zero,
  $S$ é não vazio: por exemplo, se $a \neq 0$, então $a \cdot 1 + b \cdot 0 = a$,
  e assim $|a| \in S$. Pelo princípio do bom
  ordenamento, $S$ tem um menor elemento positivo, digamos $m$, de modo que $m
  = ax_0 + by_0$ para alguns inteiros $x_0$ e $y_0$.

  Agora, provamos que $m = d$. Como $d = \text{gcd}(a, b)$, $d$ divide $a$ e
  $b$, então $a = d k$ e $b = d l$ para alguns inteiros $k$ e $l$. Substituindo
  na expressão de $m$, temos:
  \[
  m = ax_0 + by_0 = (d k) x_0 + (d l) y_0 = d (k x_0 + l y_0)
  \]
  Portanto, $m$ é um múltiplo de $d$, ou seja, $d \mid m$. Isso implica $m \geq
  d$, já que $m$ e $d$ são positivos.

  Para mostrar que $m = d$, suponha por contradição que $m > d$. Divida $a$ por
  $m$ usando o algoritmo de divisão: $a = m q + r$, onde $0 \leq r < m$. Como
  $m = ax_0 + by_0$, substituímos:
  \[
  a = (ax_0 + by_0) q + r
  \]
  Rearranjando, temos:
  \[
  r = a - (ax_0 + by_0) q = a (1 - x_0 q) + b (-y_0 q)
  \]
  Assim, $r$ é uma combinação linear de $a$ e $b$. Além disso, $0 \leq r < m$.
  Se $r > 0$, então $r \in S$, mas $r < m$, o que contradiz o fato de $m$ ser o
  menor elemento positivo de $S$. Portanto, $r = 0$, ou seja, $a = m q$, e $m$
  divide $a$. Analogamente, $m$ divide $b$. Como $m$ divide $a$ e $b$, $m$ é um
  divisor comum de $a$ e $b$, e como $d = \text{gcd}(a, b)$ é o maior divisor
  comum, temos $m \leq d$. Combinando com $m \geq d$, concluímos que $m = d$.
  Assim, $d = ax_0 + by_0$, e $d$ é uma combinação linear de $a$ e $b$.

  Na prática, $x_0$ e $y_0$ podem ser encontrados usando o Algoritmo de
  Euclides Estendido. Por exemplo, para $a = 3$, $b = 5$, aplicamos o algoritmo
  de Euclides:
  \[
  5 = 1 \cdot 3 + 2, \quad 3 = 1 \cdot 2 + 1, \quad 2 = 2 \cdot 1 + 0
  \]
  O último resto não nulo é 1, então $\text{gcd}(3, 5) = 1$. Agora, retrocedemos:
  \[
  1 = 3 - 1 \cdot 2, \quad 2 = 5 - 1 \cdot 3
  \]
   Substituímos $2$ na primeira equação:
  \[
  1 = 3 - 1 \cdot (5 - 1 \cdot 3) = 3 - 5 + 1 \cdot 3 = 2 \cdot 3 - 5
  \]
  Portanto, $1 = 2 \cdot 3 + (-1) \cdot 5$, com $x_0 = 2$, $y_0 = -1$.

  Finalmente, mostramos que qualquer combinação linear de $a$ e $b$, digamos
  $ax + by$, é um múltiplo de $d$. Como $ax + by = (ax + by) \cdot 1$, e $d$
  divide $a$ e $b$, segue que $d$ divide $ax + by$. Assim, qualquer elemento de
  $S$ é um múltiplo de $d$, e $d$ é o menor inteiro positivo em $S$.

  Este resultado é fundamental no exemplo de $\mathbb{Z}/p\mathbb{Z}$, onde $p$
  é primo e $a \in \{1, 2, \ldots, p-1\}$. Como $\text{gcd}(a, p) = 1$, o
  Teorema de Bézout garante a existência de $x$ e $y$ tais que $ax + py = 1$,
  ou seja, $ax \equiv 1 \pmod{p}$, mostrando que $a$ tem um inverso
  multiplicativo em $\mathbb{Z}/p\mathbb{Z}$, o que é essencial para provar que
  $\mathbb{Z}/p\mathbb{Z}$ é um corpo.
\end{proof}



